% from ega4-15-fr.tex
\oldpage[IV-3]{232}
on ait $g_1(z_1)=y$ et $g_1(z'_1)=y'$ (\textbf{II}, 7.1.9)~; on peut donc déjà supposer que $Y$ est le spectre d'un anneau de valuation discrète $A$, $y$ son point fermé et $y'$ son point générique. Pour tout $x\in f^{-1}(y)$, $\mathbf{k}(x)$ est une extension finie de $\mathbf{k}(y)$ par hypothèse (\textbf{II}, 6.2.2)~; on déduit de (4.5.11) qu'il existe une extension algébrique finie $k$ de $\mathbf{k}(y)$ telle que le nombre de points de $f^{-1}(y)\otimes_{\mathbf{k}(y)}k$ soit égal à $n(y)$. Comme il y a un anneau de valuation discrète $A_2$ dominant $A$ et dont le corps résiduel est fini sur $\mathbf{k}(y)$ et contient $k$ (\textbf{II}, 7.1.2), on peut en second lieu supposer que $n(y)$ est le nombre de points de $f^{-1}(y)$. Enfin, le même raisonnement montre qu'il y a une extension finie $K$ de $\mathbf{k}(y')$ telle que le nombre de points de $f^{-1}(y')\otimes_{\mathbf{k}(y')}K$ soit égal à $n(y')$~; si $w$ est une valuation de $\mathbf{k}(y')$ associée à $A$, il y a une valuation discrète de $K$ prolongeant $w$, et l'anneau $A_3$ de cette valuation domine $A$, a $K$ pour corps des fractions, et répond aux conditions du lemme.

On peut donc désormais supposer que $Y$ est le spectre d'un anneau de valuation discrète, $y$ son point fermé, $y'$ son point générique, et que $n(y)$ et $n(y')$ sont les nombres de points de $f^{-1}(y)$ et $f^{-1}(y')$, de sorte que pour tout $x\in f^{-1}(y)$ (resp. tout $x'\in f^{-1}(y')$) on a $\mathbf{k}(x)=\mathbf{k}(y)$ (resp. $\mathbf{k}(x')=\mathbf{k}(y')$).

Désignons alors par $X_i$ les sous-préschémas réduits de $X$ ayant pour espaces sous-jacents les composantes irréductibles de $X$ ($1\leqslant i\leqslant m$). En vertu de (1.10.4), la restriction $X_i\to Y$ de $f$ à tout $X_i$ est un morphisme dominant, et comme $f^{-1}(y')$ est discret, chacune des intersections $X_i\cap f^{-1}(y')$ est réduite au point générique de $X_i$~; d'où (en désignant par $n_i(z)$ le nombre géométrique de points de $X_i\cap f^{-1}(z)$ pour tout $z\in Y$), $n_i(y')=1$ pour tout $i$, et $n(y')=\sum_{i=1}^{m}n_i(y')=m$. Par ailleurs, $f^{-1}(y)$ est réunion des $X_i\cap f^{-1}(y)$, d'où
\begin{equation}
\tag{15.5.2.3}
\label{IV.15.5.2.3-fr}
  n(y)\leqslant\sum_{i=1}^{m}n_i(y),
\end{equation}
et, pour prouver (15.5.2.1), il suffit d'établir que $n_i(y)\leqslant 1$ pour tout $i$. Autrement dit, on peut supposer $X$ intègre, et le morphisme $f$ birationnel. Mais alors, pour tout $x\in f^{-1}(y)$, on a $\mathcal O_y\subset\mathcal O_x\subset K$, où $K=\mathbf{k}(y')$ est le corps des fractions de $\mathcal O_y$. Comme $\mathcal O_y$ est un anneau de valuation et que $\mathcal O_x$ domine $\mathcal O_y$, on a nécessairement $\mathcal O_y=\mathcal O_x$ (\textbf{II}, 7.1.1). Il existe alors une $Y$-section $h:Y\to X$ telle que $h(y)=x$ (\textbf{I}, 2.4.4)~; comme $Y$ est réduit et que $X$ est séparé sur $Y$, une telle section est unique (\textbf{I}, 7.2.2), donc l'ensemble $f^{-1}(y)$ ne contient qu'un seul point, ce qui achève de prouver (i).

(ii) Comme au début de (i), on est ramené à montrer que si les deux membres de (15.5.2.1) sont égaux pour toute générisation $y'$ de $y$, $f$ est propre au point $y$. Grâce au critère de propreté locale (15.7.5) et aux remarques du début on peut encore supposer que $Y=\operatorname{Spec}(A)$ est le spectre d'un anneau de valuation discrète $A$, $y$ son point fermé et $y'$ son point générique, et il s'agit alors de montrer que $f$ est propre. Utilisons (15.5.2.2) et notons que si $A'$ est un anneau de valuation discrète dominant $A$, $A'$ est un $A$-module sans torsion, donc le morphisme $\operatorname{Spec}(A')\to\operatorname{Spec}(A)$ est fidèlement plat et quasi-compact, et il revient donc au même de dire que $f$ est propre ou que le morphisme déduit de $f$ par le changement de base $\operatorname{Spec}(A')\to\operatorname{Spec}(A)$ est
